\clearpage
\chapter*{РЕФЕРАТ}
\addcontentsline{toc}{chapter}{РЕФЕРАТ}
Выпускная квалификационная работа содержит 65~с., 1~табл., 20~источников, 2~приложения.

\textbf{КЛЮЧЕВЫЕ СЛОВА:} ВОДА, ВОДЯНОЙ ПАР, ТЕРМОДИНАМИЧЕСКИЕ СВОЙСТВА, IAPWS-IF97, ЧИСЛЕННЫЕ МЕТОДЫ, МЕТОД СЕКУЩИХ, RUST, FLTK, WEBASSEMBLY, TAURI, GRPC, DOCKER, KUBERNETES, CI/CD.

Выпускная квалификационная работа посвящена разработке вычислительной модели и программного комплекса для расчёта термодинамических свойств воды и водяного пара.
В качестве нормативной основы используется промышленный стандарт \texttt{IAPWS-IF97} в редакции 2012 года~\cite{iapws_if97_2012}.
Цель работы --- разработать единое вычислительное ядро и набор каналов доставки результатов,
обеспечивающих как интерактивные инженерные расчёты,
так и интеграцию в сервисную инфраструктуру.
Задачи работы включают реализацию вычислительного ядра на языке \texttt{Rust} по стандарту \texttt{IAPWS-IF97},
поддержку прямых и обратных маршрутов расчёта
и численное расширение маршрута $(\rho, T)$
на всей рабочей области на основе метода секущих.
Также выполнены разработка настольного приложения на \texttt{FLTK},
настольной оболочки на базе \texttt{Yew}, \texttt{WebAssembly} и \texttt{Tauri},
\texttt{gRPC}-сервиса для пакетных расчётов,
контейнеризация и подготовка инфраструктуры развёртывания в \texttt{Kubernetes},
а также автоматизация сборки, тестирования и публикации артефактов
средствами непрерывной интеграции и доставки.
Объект исследования --- термодинамические свойства воды и водяного пара.
Предмет исследования --- алгоритмы и программная реализация вычисления свойств в области применимости IAPWS-IF97, а также архитектурные решения, обеспечивающие повторное использование вычислительного ядра в различных приложениях.
Методы исследования --- аналитические зависимости \texttt{IAPWS-IF97},
численные методы решения нелинейных уравнений (метод секущих),
модульное проектирование,
валидация диапазонов входных параметров
и тестирование по контрольным данным.
Результаты работы --- реализовано вычислительное ядро \texttt{if97\_core}
с единым \texttt{API},
типобезопасными единицами измерения
и унифицированной структурой результата.
Разработаны два пользовательских приложения:
нативное приложение на \texttt{FLTK}
и настольная оболочка на базе \texttt{Yew}, \texttt{WebAssembly} и \texttt{Tauri}.
Реализован \texttt{gRPC}-сервис
с разделением нагрузки ввода-вывода и вычислительной нагрузки,
а также с механизмами ограничения нагрузки.
Подготовлены \texttt{Docker}-образы,
\texttt{Kubernetes}-манифесты
и сценарии \texttt{CI/CD}.
Практическая значимость работы заключается в возможности использовать
единый набор вычислений
в инженерном настольном приложении,
в составе кроссплатформенной настольной оболочки
и как сервис для пакетных расчётов и интеграции в другие системы.
Работа состоит из введения, пяти глав, заключения, списка использованных источников и приложений.
